% !TEX root = ../main.tex
\section{Conclusion}
\label{sec:conclusion}

Across 496 runs of 4 commercial coding agents on 62 real feature-implementation tasks, only 1.6\% reach a merge-ready verdict and the entire Kubernetes track returns 0/32. The strongest harnesses pass 2.4\%, the weaker ones less than 1\%. Whatever ``ready for autonomous feature delivery'' means at Huawei's scale, the data say we are not there yet.

What we \emph{do} take away is more useful than a yes/no. The deployment-relevant gap is harness-engineering, not model scale: long prompts, scope checks, and a multi-judge gate close most of the addressable failure share, while semantic and cross-module errors remain. The audit therefore translates into a concrete rollout policy for our team---route bounded operator-tier tasks through agents under spec/scope/test gates, keep cross-cutting platform work under human authorship, and continue running the same four-judge panel on every agent draft before it reaches a reviewer. The CFP-friendly version is shorter still: \emph{the contribution is the routing rule and the gates, not a new score on a new benchmark.}

\paragraph{Future work.} Two threads are open. First, the engineering-addressable subset (coverage gaps, missing-test failures, design mismatches) is large enough to motivate harness ablations: which gates actually move the needle on real feature PRs? Second, the CapBench OSS arm makes it cheap for other teams to repeat this audit on their own agents, prompts, and review criteria; the artefact ships with that re-run path documented end-to-end.